\scanpage{27}
% loc 1–4, Soit donc A
\deleted{\unsure{Soit|Si}}Soit donc $A$ une coalgèbre (ass, \unsure{unit|unit.}) et $M$ un comodule \inserted{à g.} sur la coalgèbre. $\check A=\unsure{V|\mathbb V|W}(A)(k)$ opère à droite sur $V(M)$, donc à gauche sur $\unsure{W|V}(M)$, ou ce qui revient au $=$, sur $M$, donc \deleted{\unsure{$W($}}$M$) devient une $\check A$-\deleted{\unsure{alg}}module. Si
% loc 5–7, xi and its action
\[
\begin{aligned}
 &\xi:A\longrightarrow k\qquad(k\text{-}\mathrm{lin})\qquad\in\check A,\ x\unsure{\in|\vdash}M\\
 &\text{on a}\qquad \xi\,.\,x=(\mathrm{id}_{\unsure{M|\mu}}\otimes\xi)(u_0(x))\\
 &\text{i.e.}\qquad \xi_M=(\mathrm{id}_M\otimes\xi)\circ u_0
\end{aligned}
\]
% loc 8–10, Notons and two Hom identifications
Notons que si $\deleted{\unsure{M}}A$ est \deleted{\illegible{1} \unsure{,} \unsure{sur tout|sur tous}}\inserted{projectif de t.f. sur $k$}, on a un isom
\[
\begin{aligned}
 &\deleted{M\otimes A\simeq\mathrm{Hom}\,\mathrm{cont}_k(}\\
 &M\otimes A\simeq\mathrm{Hom}_k(\check A,M)
\end{aligned}
\]
% loc 11–13, bracketed u0 and equivalence
donc la donnée de
\tikz[baseline=(h.base)]{\node[inner sep=2pt](h){$M\overset{u_0}{\longrightarrow}M\otimes_k A$};\draw (h.south west) .. controls +(-.3,0) and +(-.12,-.15) .. (h.north west) -- (h.north east) .. controls +(-.025,-.2) and +(-.1,-.35) .. (h.south east);}
[satisfaisant certaines conditions) \deleted{\unsure{revient}}\inserted{\uline{équivaut}} à la donnée de
\[
 \check A\longrightarrow\mathrm{Hom}_k(M,M)
\]
% loc 14–18, satisfying conditions, Le = argument
satisfaisant aux conditions précédentes. On trouve alors que la donnée d'une structure de \inserted{$A$-}comodule \deleted{\unsure{équivaut à}} sur $M$. Le $=$ argument nous montre que sur le $k$-module $A$ projectif de t.f. $A$, la donnée d'une structure de coalgèbre (ass. \unsure{unit.|un?.}) équivaut à celle d'une str. d'algèbre (ass. unit.) sur $\check A$.

% loc 19–20, cancelled starts and interlinear replacements
\deleted{Si $k$ est un corps}, On a un \unsure{rsult.|résult.|résultat} analogue \deleted{\unsure{sans supposer}}\inserted{si $A$ est libre}, \deleted{\unsure{$A$ de t.f., alors}} \inserted{p.ex. $k$ un corps} ; plus 
% loc 20–23, left margin and struck limit argument
% slanted
\margin{gén. si $A$ est somme \unsure{directe} de modules \unsure{proj.} de t.f.}
\deleted{\unsure{$A$ de t.f.} En effet, \unsure{on aura} $A=\underset{\longrightarrow}{\lim}_{\unsure{i|1}}\unsure{A_i|A'}$, \unsure{$A_i$|$A'$} les sous-modules de dim finie, donc $\check A=\underset{\longleftarrow}{\lim}\unsure{A_i|A_i'|A'}$ (\unsure{lim. proj. de} \unsure{duals|duels|produits} de \unsure{vectoriels|vecteurs} de dim finie) de $\check A$}
% loc 23–26, convergence and Hom isomorphisms
Munissant $\check A$ de la topologie de la convergence simple, on aura
\[
 M\otimes_k A\simeq\mathrm{Hom}\,\mathrm{cont}_k(\check A,M)
\]
donc
\[
 \mathrm{Hom}(\unsure{N|W},M\otimes_{\unsure{k|u}}A)
 \simeq\mathrm{Hom}\,\mathrm{cont}_{\unsure{K|k}}\bigl(\check A,\,
 \underbrace{\deleted{\unsure{H}}\mathrm{Hom}_k(N,M)}_{\substack{\text{muni de}\\\text{la top. de la}\\\text{conv. simple}}}\bigr)
\]
% loc 27–28, sentence continuing on scan 28
donc on trouve qu'une structure de \unsure{comodule} (\unsure{ass.} unit) sur $M$ revient à la donnée d'un
% archive sheet number: 26
